\documentclass[11pt]{article}
\usepackage[T1]{fontenc}
\usepackage[utf8]{inputenc}
\usepackage{lmodern}
\usepackage[margin=1in]{geometry}
\usepackage{amsmath,amssymb,amsthm,mathtools}
\usepackage{hyperref}
\usepackage{enumitem}
\usepackage{verbatim}
\newcommand{\OPHPaperReleaseID}{r1381}
\newcommand{\OPHPaperReleaseDate}{April 10, 2026}
\newcommand{\OPHPaperReleaseBanner}{%
\begin{center}
\small\textbf{Paper release:} \texttt{\OPHPaperReleaseID}\quad\textbf{Released:} \OPHPaperReleaseDate
\end{center}
}


\hypersetup{colorlinks=true,linkcolor=blue,citecolor=blue,urlcolor=blue}

\title{Screen Microphysics, Patches, and Observer Synchronization in OPH}
\author{B.\ M\"uller}
\date{\OPHPaperReleaseDate}

\begin{document}
\maketitle
\OPHPaperReleaseBanner

\begin{abstract}
This paper is a phase-1 architecture note for Observer-Patch Holography (OPH). The
objective is constructive and simulator-facing: choose one finite local screen model that can
actually be encoded in qubits, make the map from local registers to patch algebras and overlap
observables explicit, and state how stable records, observer patterns, local interaction, and
observer-level synchronization could be built inside that same microscopic system. The working
choice is a constrained finite gauge-register or quantum-link style model on a cellulated
screen, used here as a \emph{reference architecture} rather than as a claim that the unique OPH
UV completion has already been identified. We separate reused OPH structure, architecture-level
claims, simulator targets, and conjectural bridges. In particular, this note does \emph{not}
claim that the OPH continuum or gravity branch has already been derived from the present
microphysics. The deliverable of this round is a sharper blueprint: local Hilbert spaces,
qudit-to-qubit encoding agenda, patch and overlap observable dictionary, observer criterion,
synchronization and repair options, a first implementable reference model, and a validation
package that can fail.
\end{abstract}

\tableofcontents
\newpage

\section{Project Charter}

\textbf{Status.} Extra side paper. Not part of the compact recovered core. Not yet part of the
release-critical OPH theorem stack.

\textbf{Purpose.}
\begin{enumerate}[leftmargin=2em]
\item Give a concrete reference architecture for local finite screen degrees of freedom.
\item Make the bridge from local registers to patches, overlaps, records, observers, and
synchronization explicit enough to build toy simulators.
\item Keep theorem claims, simulator claims, and conjectural bridges sharply separated.
\item Feed the existing OPH microphysics, repair, and observer task lanes without pretending
to close them automatically.
\end{enumerate}

\textbf{Non-goals.}
\begin{enumerate}[leftmargin=2em]
\item Do not claim that the unique UV completion is known.
\item Do not claim that the OPH scaling-limit continuum or gravity branch is already derived
here.
\item Do not treat ``observer'' as a consciousness claim.
\item Do not treat synchronization of overlap data as a proof of full continuum closure.
\end{enumerate}

\textbf{Phase-1 output.} A reference microscopic screen model, a qubit-encoding agenda, a patch
and overlap dictionary, an observer criterion, a synchronization design space, and a simulator
validation package.

\section{Claim Taxonomy}

The point of this paper is not to maximize the number of statements called ``derived.'' The
point is to keep the layers clean.

\begin{enumerate}[leftmargin=2em,label=\textbf{T\arabic*.}]
\item \textbf{Reused OPH structure.} The OPH screen-net picture, overlap consistency language,
local MaxEnt and refinement-stability branch, recoverable generalized entropy package, and the
regulator-scale edge-center completion language are reused as background structure. This paper
does not re-prove them.

\item \textbf{Architecture claims.} Once one finite gauge-register reference architecture is
chosen, the patch algebra, overlap algebra, edge-sector observables, record layer, and
observer-level interfaces can be defined explicitly at fixed cutoff.

\item \textbf{Simulator claims.} Small-lattice digital or tensor-network simulations can test
whether the reference architecture exhibits gauge-invariant patch observables, stable records,
low conditional mutual information across collars, and robust synchronization under local repair
schedules.

\item \textbf{Conjectural bridge claims.} It is plausible that an appropriate refinement-stable
phase of such models lands in the OPH modular-geometric branch and later supports the OPH
continuum or gravity story, but that bridge is not proved here.

\item \textbf{Excluded claims.} This note does not identify the final OPH microphysics, does not
establish a unique repair law, and does not close the Standard Model or continuum branches.
\end{enumerate}

\section{Why a Finite Gauge-Register Reference Architecture}

Several microscopic families are logically possible: constrained finite gauge registers,
quantum-link truncations of compact gauge systems, string-net or quantum-double models,
generic tensor-network ans\"atze, or ad hoc observer automata attached to a screen. For a
phase-1 constructive draft, the first working model should be selected by engineering
criteria rather than by aesthetic breadth.

The selection criteria are:
\begin{enumerate}[leftmargin=2em]
\item finite local Hilbert spaces at fixed cutoff;
\item exact local constraints rather than only emergent ones;
\item a natural patch and overlap algebra;
\item explicit edge-sector or center observables on collars;
\item straightforward qudit-to-qubit compilation;
\item a path from local mismatch syndromes to repair or synchronization maps;
\item extensibility from toy groups such as $\mathbb Z_2$ or $S_3$ to richer quantum-link
truncations.
\end{enumerate}

A constrained finite gauge-register or quantum-link style screen model is therefore the
default reference architecture for this draft. It is finite-dimensional, local, gauge-aware,
and already speaks the same language as the OPH regulator package. It also makes the collar
center and edge-sector story concrete instead of metaphorical.

This choice is a working \emph{reference architecture}. It is not a claim that all other
microscopic families are ruled out.

\section{Reference Architecture: Local Screen Registers on a Cellulated Sphere}

Take a finite cellulation $\Gamma=(V,E,F)$ of a horizon screen $S^2$. Each oriented link
$e\in E$ carries a finite register $\mathcal H_e$. The simplest version is a finite-group
register with basis $\{|g\rangle_e:g\in G\}$ for a finite group $G$, so
\begin{equation}
\mathcal H_e \cong \mathbb C^{d_G},\qquad d_G=|G|.
\end{equation}
A more flexible version replaces $|g\rangle$ by a truncated quantum-link basis labeled by a
small irrep set and matrix indices. For phase 1, the finite-group presentation is easier to
simulate, while the quantum-link presentation shows how the same pattern could extend to
compact-group truncations.

To support records and observer logic, attach a small record register to selected vertices or
coarse cells:
\begin{equation}
\mathcal H_v^{\mathrm{rec}} \cong (\mathbb C^2)^{\otimes m_v}.
\end{equation}
The extended Hilbert space is then
\begin{equation}
\widetilde{\mathcal H}_{\Gamma}=
\bigotimes_{e\in E}\mathcal H_e
\otimes
\bigotimes_{v\in V}\mathcal H_v^{\mathrm{rec}}.
\end{equation}

At each vertex $v$, a local gauge action acts on the incident oriented links by left/right
multiplication or by the corresponding truncated quantum-link action. Let $P_v$ denote the
Gauss projector onto the locally admissible subspace. The physical cutoff Hilbert space is
\begin{equation}
\mathcal H_{\mathrm{phys}} = \Bigl(\prod_{v\in V}P_v\Bigr)\widetilde{\mathcal H}_{\Gamma}.
\end{equation}

A minimal reference Hamiltonian or Trotterized circuit family should contain four layers:
\begin{equation}
H_{\mathrm{ref}} =
\lambda_G\sum_{v}(1-P_v)
+\lambda_B\sum_{f}(1-B_f)
+\lambda_R\sum_{a}(1-R_a)
+H_{\mathrm{sync}},
\end{equation}
where $B_f$ are plaquette or face operators, $R_a$ are record-stabilization or record-write
terms, and $H_{\mathrm{sync}}$ is a local overlap-repair interaction to be chosen from a small
menu later in the paper. The exact coefficients and update schedules are not fixed here.
What matters is the architecture: gauge constraints, face dynamics, record dynamics, and
repair dynamics all live on one finite screen.

\begin{verbatim}
Microscopic screen picture

   record q(v1)      record q(v2)      record q(v3)
        o------------------o------------------o
        | \              / | \              / |
        |  \    f1      /  |  \    f2      /  |
        |   \          /   |   \          /   |
        |    \        /    |    \        /    |
        o-----\------/-----o-----\------/-----o
      q(v4)    \    /    q(v5)    \    /    q(v6)
                 \  /                \/
                  \/                 /\
                  /\                /  \
                 /  \              / f4 \
                / f3 \            /      \
               /      \          /        \
              o--------\--------o----------o
          record q(v7)          q(v8)     q(v9)

   links carry finite gauge registers
   vertices host Gauss checks and optional record qubits
   faces carry plaquette/holonomy dynamics
\end{verbatim}

The simulator-facing question is therefore not ``what is the final UV completion?'' It is
``can one finite local screen model realize the OPH patch, edge, record, and synchronization
interfaces well enough to be testable?''

\section{Qudit-to-Qubit Encoding Agenda}

If the screen model is ever going to be built or simulated on conventional hardware, the
local finite registers must be compiled into qubits. This paper does not pick a single global
encoding forever; it sets the agenda and the tradeoffs.

\subsection{Binary, unary, and factorized encodings}

For a $d$-level link register, the default binary encoding uses
\begin{equation}
 n_e = \lceil \log_2 d\rceil
\end{equation}
qubits per link, with a penalty term suppressing unused binary strings when $d$ is not a
power of two. Binary encoding is qubit-efficient and natural for digital simulation.

Unary or one-hot encoding uses $d$ qubits per link with one excitation allowed. It is less
qubit-efficient but makes local projectors and permutation actions shallower. This is useful
for very small groups and for verification runs.

For truncated quantum-link bases labeled by irreps and matrix indices, a factorized encoding
is often cleaner:
\begin{equation}
|R,m,n\rangle \longmapsto |R\rangle_{\mathrm{irrep}}\otimes |m\rangle_{\mathrm{left}}
\otimes |n\rangle_{\mathrm{right}}.
\end{equation}
This keeps the representation label explicit, which is attractive when collar sector labels
are later measured.

\subsection{Agenda for phase 1}

The encoding work should proceed in the following order.
\begin{enumerate}[leftmargin=2em]
\item Start with $G=\mathbb Z_2$ or $\mathbb Z_N$ in binary encoding to validate the patch,
overlap, and synchronization pipeline cheaply.
\item Move to $S_3$ or another small nonabelian finite group to validate nonabelian sector
projectors, collar labels, and character-based observables.
\item Add a truncated quantum-link variant only after the finite-group version produces stable
measurements and sensible scaling diagnostics.
\item Compare binary and unary encodings only at the level of compilation cost and numerical
stability, not as a statement about final physics.
\end{enumerate}

What must be compiled into qubit circuits is now explicit:
\begin{enumerate}[leftmargin=2em]
\item local gauge transformations and Gauss projectors;
\item plaquette or loop operators;
\item record-write and record-read channels;
\item overlap projectors, edge-sector projectors, and mismatch syndromes;
\item local repair maps that preserve gauge admissibility.
\end{enumerate}

\section{From Local Registers to Patches and Overlaps}

The bridge from local finite registers to OPH patch language has to be explicit. A patch is
not only a geometric region. It is a connected region together with the gauge-invariant
observable algebra available there.

Let $P\subseteq \Gamma$ be a connected collection of faces, together with all vertices and
links incident to those faces. Its extended cutoff Hilbert space is
\begin{equation}
\widetilde{\mathcal H}_P=
\bigotimes_{e\subset P}\mathcal H_e
\otimes
\bigotimes_{v\subset P}\mathcal H_v^{\mathrm{rec}},
\end{equation}
and its extended local algebra is
\begin{equation}
\widetilde{\mathcal A}(P)=\mathcal B(\widetilde{\mathcal H}_P).
\end{equation}
The physical patch algebra is the boundary-fixed subalgebra
\begin{equation}
\mathcal A(P)=\widetilde{\mathcal A}(P)^{G_{\partial P}},
\end{equation}
in the same spirit as the regulator picture already used in OPH.

If $P$ and $Q$ overlap, the shared region $B=P\cap Q$ is not merely a geometric set. It is
where the shared or comparable observables live. At fixed cutoff, the collar Hilbert space is
expected to decompose as
\begin{equation}
\mathcal H_B \cong \bigoplus_{\alpha}
\bigl(H_{b_L^{\alpha}}\otimes H_{b_R^{\alpha}}\bigr),
\end{equation}
with center projectors $\Pi_{\alpha}$ labeling edge sectors. That decomposition is not proved
again here; it is reused OPH regulator language and serves as the observable interface for the
reference model.

\subsection{Patch and overlap observable dictionary}

The working dictionary is the core of this paper.

\begin{enumerate}[leftmargin=2em]
\item \textbf{Link register basis / electric occupation.}
The local computational basis on a link is the microscopic screen data. In a finite-group
presentation this is the basis $|g\rangle_e$; in a quantum-link truncation it is the local
representation basis.

\item \textbf{Gauss projector $P_v$.}
This is the local admissibility check. It says which neighboring link configurations may be
regarded as physical around a vertex.

\item \textbf{Plaquette operator $B_f$ or loop operator $W_\gamma$.}
This is the patch-level probe of curvature, flux, or sector excitation. Closed loops contained
inside a patch belong to $\mathcal A(P)$.

\item \textbf{Boundary or collar sector projectors $\Pi_{\alpha}$.}
These are the overlap-accessible center observables. They are the natural candidates for
shareable edge data and for robust cross-patch records.

\item \textbf{Record observables $Z^{\mathrm{rec}}_{v,a}$ and pointer subalgebras.}
These are commuting or approximately commuting observables stabilized long enough to act as
memories. A record is useful only if it can be read repeatedly with low back-action on the
shared sector data.

\item \textbf{Patch state $\rho_P$.}
The reduced state on $\mathcal A(P)$ is the patch description accessible to an internal
pattern supported on $P$.

\item \textbf{Overlap mismatch syndrome $S_{PQ}$.}
This is a function of shared observables, typically of the collar sector labels and the record
comparisons on $B=P\cap Q$. It is the direct input to any repair or synchronization update.
\end{enumerate}

\begin{verbatim}
Patch/overlap interface

   Patch P bulk             shared collar B            Patch Q bulk
  ----------------         =================         ----------------
  | loops W_gamma |        || Pi_alpha     ||        | loops W_gamma |
  | records R_P   | <----> || shared bits  || <----> | records R_Q   |
  | local state   |        || syndrome S   ||        | local state   |
  ----------------         =================         ----------------

   what is compared on B is not the whole quantum state
   it is a selected shared observable family on the overlap
\end{verbatim}

This is where the bridge to observers becomes concrete. An observer never needs access to the
full global state. It needs persistent access to some patch algebra and to a stable shared
observable family on overlaps with neighboring patches.

\section{Records and the Observer Criterion}

The word ``observer'' should be used operationally. In OPH language, an observer is an
internal pattern with patch access and records. In a finite screen model, that can be stated
more sharply.

\subsection{Records}

A \emph{record layer} is a commuting or approximately commuting subalgebra of robust pointer
observables, supported partly inside a patch and partly on the boundary sector interface. A
useful record variable must satisfy three practical conditions:
\begin{enumerate}[leftmargin=2em]
\item it is written by local interactions with nearby screen registers;
\item it survives for many local update cycles compared with the readout timescale;
\item it can be compared across overlaps through shared observables or recoverable summaries.
\end{enumerate}

This prevents the word ``record'' from collapsing into an arbitrary entangled correlation.
The role of the record layer is to produce stable, repeatedly accessible data.

\subsection{Working observer criterion}

The OPH manuscript already uses the tuple
\begin{equation}
O=(P_O,\mathcal A(P_O),\rho_O,R_O)
\end{equation}
for an observer patch, local algebra, local state, and records. For the reference
microphysics we add an explicit update interface and treat an observer as a pattern class
rather than as a primitive object:
\begin{equation}
O_{\mathrm{ref}}=(P_O,\mathcal A(P_O),\rho_O,R_O,\mathcal U_O).
\end{equation}
Here $\mathcal U_O$ is a local channel or update family that writes, compares, and repairs
records using only observables available in $P_O$ and on its overlaps.

A subsystem or pattern counts as an \emph{observer in the present operational sense} if it
satisfies all of the following.
\begin{enumerate}[leftmargin=2em]
\item \textbf{Patch access.} It is supported on a connected patch $P_O$ with a nontrivial
shared overlap algebra to neighboring patches.
\item \textbf{Record support.} It carries a set $R_O$ of metastable record observables.
\item \textbf{Update capability.} It can condition future local actions on those records.
\item \textbf{Comparison capability.} It can query at least one shared overlap observable
family and detect mismatch syndromes with neighbors.
\item \textbf{Persistence.} The pattern remains identifiable for many local update cycles.
\end{enumerate}

This is an observer pattern criterion, not a consciousness criterion. It is meant to support
simulator definitions of record-bearing internal agents or agent-like subsystems.

\begin{verbatim}
Local registers --> records --> observer pattern --> synchronization

   link qudits / qubits
          |
          v
   gauge-invariant patch observables
          |
          v
   robust pointer variables / records
          |
          v
   persistent update-capable pattern O_ref
          |
          v
   compare shared overlap data and repair mismatch
\end{verbatim}

\section{Interaction, Synchronization, and Repair}

The main missing bridge in many holographic narratives is the move from ``overlaps exist'' to
``observers actually synchronize.'' For the present reference architecture that bridge should be
an explicit local interface, not a metaphor. A synchronization step is a local update on one
overlap $B=P\cap Q$ that compares a declared family of shared observables, produces a finite
mismatch syndrome, and applies a gauge-admissible correction while preserving the record layer.

\subsection{What is synchronized}

Full quantum states are not synchronized. The synchronization target is a chosen overlap
observable family,
\begin{equation}
\mathcal O^{\mathrm{sh}}_{PQ}
=
\{\Pi_{\alpha}^{(PQ)}\}_{\alpha}
\cup
\{Q_a^{(PQ)}\}_a
\cup
\{M_r^{(PQ)}\}_r
\subseteq
\mathcal A(P\cap Q),
\end{equation}
typically built from collar sector projectors, gauge-invariant boundary charges, and read-stable
record summaries. In practical terms, synchronization means that two observers come to agree on
\emph{the values or distributions of shared observables}, not that their internal microscopic
states become identical. For a phase-1 build it is enough that each overlap expose one collar
sector label, one small set of gauge-invariant boundary observables, and one record summary.

A convenient mismatch syndrome is
\begin{equation}
S_{PQ}
=
\left(
1-\delta_{\alpha_P,\alpha_Q},
\{Q_a(P)-Q_a(Q)\}_a,
\{M_r(P)-M_r(Q)\}_r
\right),
\end{equation}
where $\alpha_P$ and $\alpha_Q$ are the collar-sector labels inferred from the two patch sides.
The associated local mismatch cost may be taken to be
\begin{equation}
\Xi_{PQ}
=
\lambda_{\alpha}(1-\delta_{\alpha_P,\alpha_Q})
+
\sum_a \lambda_a |Q_a(P)-Q_a(Q)|^2
+
\sum_r \lambda_r |M_r(P)-M_r(Q)|^2,
\end{equation}
with nonnegative design weights. The global mismatch potential is then
\begin{equation}
\Phi = \sum_{\langle P,Q\rangle} w_{PQ}\,\Xi_{PQ}.
\end{equation}
Nothing in this paper proves that a chosen repair law must decrease $\Phi$ under all schedules.
The point is narrower: the microscopic interface on which such claims would later be tested is
now explicit.

\subsection{Phase-1 repair interface}

A single local repair step on the overlap between $P$ and $Q$ should have schematic form
\begin{equation}
\mathcal R_{PQ}
=
\mathcal W_{PQ}\circ \mathcal C_{PQ}\circ \mathcal M_{PQ}\circ \mathcal P_{PQ},
\end{equation}
with four logically separate pieces.
\begin{enumerate}[leftmargin=2em]
\item $\mathcal P_{PQ}$ is a local admissibility check or projection that verifies the Gauss
constraints on the support of the planned update.
\item $\mathcal M_{PQ}$ extracts the chosen mismatch syndrome $S_{PQ}$ from the shared
observable family $\mathcal O^{\mathrm{sh}}_{PQ}$.
\item $\mathcal C_{PQ}$ selects one correction from a finite local menu
$\{C_{PQ}^{(k)}\}_k$ supported on the overlap collar and its immediate neighborhood.
\item $\mathcal W_{PQ}$ rewrites or refreshes the declared record summaries after the correction
has been accepted.
\end{enumerate}

For a simulator-facing build the correction menu should be intentionally small. A useful first
menu is: flip one boundary-adjacent link register, apply one local face move, refresh one stale
record bit, or do nothing. Accepted corrections should satisfy two practical tests:
\begin{enumerate}[leftmargin=2em]
\item \textbf{Gauge admissibility.} The update preserves or restores the local Gauss
constraints on its support.
\item \textbf{Monotone local fit.} The update strictly lowers $\Xi_{PQ}$, or leaves it
unchanged while decreasing a declared tie-breaker such as record disagreement count or support
size.
\end{enumerate}

This is enough to define a repair loop without pretending that termination, confluence, or
continuum robustness have already been proved.

\subsection{Repair mechanism families}

Three local mechanism families are worth keeping on the table.

\paragraph{Option A: syndrome-based overlap repair.}
Measure or infer a mismatch syndrome on $P\cap Q$, apply a gauge-covariant correction on one or
both sides, and keep the update local. In the present note this means an explicit correction
menu $\{C_{PQ}^{(k)}\}$ acting on the same overlap observables that define $S_{PQ}$. This is
closest to the distributed-consensus language of the separate consensus paper.

\paragraph{Option B: dissipative synchronization.}
Use local Lindbladian or CPTP channels that preserve the Gauss constraints while damping
mismatch in the shared observable family. This is attractive for noisy simulation because it
builds repair into the dynamics, but it should be treated as a second-round upgrade rather than
the first debugging target.

\paragraph{Option C: record-mediated Bayesian or modular update.}
When a robust record layer already exists, synchronize the record summaries first and only then
feed the result back into local control choices. This is the most observer-facing option but is
also the furthest from a bare microscopic implementation.

\subsection{Recommended phase-1 choice}

The first implementable build should use Option A as the default synchronization mechanism: a
local syndrome-reduction map on overlaps, with gauge checks before and after the correction.
For a discrete asynchronous schedule, one update cycle should be:
\begin{enumerate}[leftmargin=2em]
\item choose one neighboring patch pair $(P,Q)$ from the patch-overlap graph;
\item measure or infer the current syndrome $S_{PQ}$;
\item if $S_{PQ}=0$, leave the overlap untouched and log a no-op;
\item otherwise choose one correction $C_{PQ}^{(k)}$ from the finite local menu;
\item recheck local gauge admissibility and reject the move if it fails;
\item refresh the local record summary observables and log the new $\Xi_{PQ}$ and $\Phi$.
\end{enumerate}

This choice has four advantages.
\begin{enumerate}[leftmargin=2em]
\item It is nearest to the existing OPH repair vocabulary.
\item It turns disagreement into a measurable object.
\item It can be tested on small finite graphs with asynchronous schedules.
\item It leaves open the later upgrade to dissipative or modular versions.
\end{enumerate}

What this paper can responsibly say is the following: if a chosen repair family is local,
gauge-covariant, terminating, and locally confluent on the selected shared observable family,
then observer-level synchronization is schedule-robust on that family. The proof of that kind
of statement belongs to the separate consensus paper. The present task is to supply the
microphysical inputs that such repair maps act on.

\section{What a First Implementable Reference Model Should Contain}

A first actual build should be deliberately small, explicit, and reproducible. The right first
target is not ``the final UV completion'' but one finite regulated screen on which every object
named earlier in the paper is actually instantiated.

\subsection{Concrete baseline: octahedral $\mathbb Z_2$ pilot}

Take the octahedral cellulation of $S^2$ and put one qubit on each oriented link with
computational basis $\{|0\rangle_e,|1\rangle_e\}$. A minimal gauge layer may then be written as
\begin{equation}
A_v=\prod_{e\ni v} X_e,
\qquad
P_v=\frac{1+A_v}{2},
\qquad
B_f=\prod_{e\subset \partial f} Z_e,
\end{equation}
where $P_v$ is the local Gauss projector and $B_f$ is the plaquette flux operator. This is the
smallest useful debugging environment because the physical subspace, plaquette sectors, and
overlap parities can all be checked by hand on a small graph.

Choose three overlapping cap patches $P_1,P_2,P_3$ so that each pair shares a nontrivial collar.
Attach one record qubit $r_i$ to each patch and declare $Z_{r_i}$ to be the patch-level record
observable. On each overlap declare the shared observable family
\begin{equation}
\mathcal O^{\mathrm{sh}}_{P_iP_j}
=
\{\Pi_{+}^{(ij)},\Pi_{-}^{(ij)},Q_{\partial}^{(ij)},Z_{r_i}Z_{r_j}\},
\end{equation}
where $\Pi_{\pm}^{(ij)}$ project onto even/odd collar-flux parity and
$Q_{\partial}^{(ij)}$ is one chosen gauge-invariant boundary observable for that collar.

The corresponding phase-1 update loop should be equally explicit:
\begin{enumerate}[leftmargin=2em]
\item apply the Gauss-check layer;
\item apply one plaquette layer on the chosen local face set;
\item write or refresh the record qubits from nearby gauge-invariant observables;
\item choose one overlap $(P_i,P_j)$ and evaluate $S_{P_iP_j}$;
\item if needed, apply one local correction from a finite lookup table;
\item reproject to the admissible subspace and log the resulting observables.
\end{enumerate}

This already produces a concrete simulator object: an octahedral $\mathbb Z_2$ screen with
patches, overlap observables, record bits, and an asynchronous repair law.

\subsection{First nonabelian upgrade: $S_3$}

Once the $\mathbb Z_2$ pilot is stable, the first nonabelian upgrade should replace each link
qubit by a six-state $S_3$ register or an encoded qubit representation thereof. The same patch
cover and repair loop can then be reused, but the collar observables should now include class or
irrep-sensitive sector labels rather than only parity data. The point of the $S_3$ step is not
numerical scale; it is to test whether the record and synchronization interface still behaves
well when the edge sector is genuinely nonabelian.

\subsection{Minimal ingredient checklist}

The first reference build should therefore contain the following ingredients.
\begin{enumerate}[leftmargin=2em]
\item \textbf{A small spherical cellulation.} Octahedral first, icosahedral only after the pilot
works.
\item \textbf{One explicit gauge family.} $\mathbb Z_2$ for the first complete debug loop, $S_3$
for the first nonabelian test.
\item \textbf{Explicit registers.} One encoded link register on every oriented edge and one small
record register on each declared observer patch.
\item \textbf{Explicit local layers.} Gauss check, plaquette update, record write, overlap
compare, repair, and post-update admissibility check.
\item \textbf{A coarse patch cover.} At least three overlapping patches so schedule dependence is
not hidden by a single-pair toy geometry.
\item \textbf{A declared shared observable family.} Sector projectors, one boundary observable,
and one record summary observable per overlap.
\item \textbf{A declared observer pattern.} At least one persistent record-bearing subsystem that
uses its records to choose later local actions.
\end{enumerate}

That model would still be a toy. But it would already answer the right phase-1 question: can
one finite local screen carry patches, edge sectors, records, observers, and stable local
synchronization in one consistent simulation?

\section{Simulator Validation Package}

A useful simulator package has to be able to fail. The point is not to accumulate plots; it is
to identify the smallest set of logged quantities that discriminate between a working fixed-cutoff
observer/repair architecture and a merely decorative toy model.

\subsection{Core quantities to log}

For every run the simulator should log at least the following quantities:
\begin{align}
g(t) &= \sum_v \langle 1-P_v\rangle_t, \\
\Phi(t) &= \sum_{\langle P,Q\rangle} w_{PQ}\,\Xi_{PQ}(t), \\
C_r(\Delta t) &= \langle Z_r^{\mathrm{rec}}(t)\,Z_r^{\mathrm{rec}}(t+\Delta t)\rangle, \\
D_{\mathrm{sched}}
&=
\max_{s,s'}
\sum_{\langle P,Q\rangle}
\|\mu_{PQ}^{(s)}-\mu_{PQ}^{(s')}\|_1.
\end{align}
Here $g(t)$ is the gauge-violation rate, $\Phi(t)$ is the aggregate overlap mismatch,
$C_r(\Delta t)$ measures record persistence, and $D_{\mathrm{sched}}$ compares the final shared
observable distributions $\mu_{PQ}^{(s)}$ obtained from different asynchronous schedules $s$.

\subsection{Kinematic validation}
\begin{enumerate}[leftmargin=2em]
\item Verify the encoded Gauss projectors and the dimension of the physical subspace on the
chosen finite graph.
\item Verify that patch algebras restricted to gauge-invariant operators behave correctly under
inclusion and overlap.
\item Verify that the chosen collar projectors really produce a measurable sector
decomposition.
\item Require that $g(t)$ stay at numerical or shot-noise level during ordinary runs; if it
does not, the repair dynamics are not even acting on the declared physical space.
\end{enumerate}

\subsection{Record and observer validation}
\begin{enumerate}[leftmargin=2em]
\item Measure record lifetimes under the microscopic dynamics.
\item Check repeated-read stability of the declared record observables.
\item Check whether records can be compared across overlaps without uncontrolled back-action on
the shared sector data.
\item Require that the declared observer pattern remain identifiable for many update cycles
relative to the readout cadence; otherwise the paper's operational observer criterion has not
been instantiated.
\end{enumerate}

\subsection{Synchronization validation}
\begin{enumerate}[leftmargin=2em]
\item Run multiple asynchronous repair schedules from the same initial family and compare the
resulting shared observable outputs.
\item Track the mismatch potential $\Phi$ and test whether it decreases under the chosen local
repair family.
\item Introduce controlled defects or frustrated cycles and determine whether the repair map
correctly distinguishes removable mismatch from topologically stable residual obstruction.
\item Treat small $D_{\mathrm{sched}}$ on the declared shared observables as the operational
signature of schedule robustness at fixed cutoff. Large $D_{\mathrm{sched}}$ means the repair law
has not yet achieved observer-level synchronization on that family.
\end{enumerate}

\subsection{Negative controls}
\begin{enumerate}[leftmargin=2em]
\item Break gauge covariance deliberately and check that patch or overlap observables become
schedule-sensitive.
\item Remove the record-stabilization layer and check that observer-level synchronization
collapses even if microscopic overlaps still exist.
\item Replace the declared shared observable family by an overambitious one and check that the
repair map fails to converge.
\end{enumerate}

\subsection{Reading rule for the validation package}

Even a successful validation package would not by itself prove the OPH scaling-limit gravity
branch, would not establish a unique UV completion, and would not close the quantitative gauge
or matter branches. It would only show that a finite local screen model can realize the needed
fixed-cutoff architecture. Conversely, failure is informative: if $g(t)$ stays high, if
$C_r(\Delta t)$ decays too quickly, if $\Phi$ does not reliably fall, or if
$D_{\mathrm{sched}}$ remains large, then that reference architecture does not yet realize the
interface this paper claims to make explicit.

\section{Dependence on Existing OPH Results and Open Bridges}

This draft depends heavily on already existing OPH structure. The point is to reuse it
cleanly and to say where the bridges remain open.

\subsection{Imported OPH structure}

The following ingredients are treated as imported OPH language or results.
\begin{enumerate}[leftmargin=2em]
\item The screen-net viewpoint on a horizon screen $S^2$ with patch algebras.
\item Overlap consistency as the requirement that local descriptions agree on shared
observables.
\item The local MaxEnt and refinement-stability branch as the current OPH dynamical baseline.
\item The recoverable generalized entropy package and the collar language in which conditional
mutual information and recovery become meaningful.
\item The regulator-scale edge-center completion picture, including the idea that a collar
supports center labels or edge sectors visible from both sides.
\item The observer tuple $(P_O,\mathcal A(P_O),\rho_O,R_O)$ as an existing OPH access model.
\end{enumerate}

\subsection{Claims specific to this note}

Relative to those imported ingredients, this note claims only a limited fixed-cutoff package.
\begin{enumerate}[leftmargin=2em]
\item One concrete finite-register reference architecture for a screen microphysics program.
\item One explicit dictionary from local registers to patches, overlaps, records, observers, and
shared observables.
\item One explicit phase-1 repair interface and update loop for overlap synchronization.
\item One concrete $\mathbb Z_2 \rightarrow S_3$ build ladder and one validation package that can
succeed or fail.
\end{enumerate}

These are architecture-level and simulator-level claims. They are not yet continuum claims, not
yet Born-rule theorems, and not yet quantitative Standard Model closures.

\subsection{Open conjectural bridges}

The following steps remain open and must not be smuggled in as completed results.
\begin{enumerate}[leftmargin=2em]
\item Showing that the chosen finite gauge-register architecture realizes the correct
refinement-stable MaxEnt phase rather than merely imitating its language.
\item Deriving the edge-sector weight law used in the compact quantitative branch, rather than
only providing finite models on which it can be measured.
\item Showing that the collar recovery and edge-sector behavior of the reference model scale in
the way required by the OPH geometric modular branch.
\item Proving termination, confluence, and gauge covariance of the repair laws in the consensus
paper, rather than only defining the microscopic objects they act on.
\item Lifting the record layer, observer criterion, and checkpoint data to theorem-level
measurement and continuation packages.
\item Deriving the OPH continuum or gravity branch from this microphysics rather than simply
using compatible notation.
\item Showing that one reference architecture is unique or even preferred beyond the current
engineering criteria.
\item Showing that the same microscopic build closes the quantitative gauge, electroweak, or
matter branches.
\end{enumerate}

The structure of this paper is therefore asymmetric on purpose: the fixed-cutoff architecture
is concrete; the scaling-limit bridges remain explicit research problems.

\section{Relation Back to the Existing Completion Tree}

This extra paper should be read as a phase-1 bridge document. It does not replace the target
task-tree papers. It makes explicit the finite-cutoff objects that those later tasks would need
to reference or prove things about.

\subsection{Direct feed: \texttt{papers.compact.e.24}}

The compact microphysics anchor needs explicit fixed-cutoff sector data rather than a purely
formal label. This note now supplies:
\begin{enumerate}[leftmargin=2em]
\item finite screen architectures with declared collar observables $\Pi_{\alpha}^{(PQ)}$;
\item explicit mismatch costs $\Xi_{PQ}$ and aggregate potential $\Phi$;
\item concrete $\mathbb Z_2$ and $S_3$ pilot models on which edge-sector weighting laws can be
measured.
\end{enumerate}
What it does \emph{not} yet supply is the derivation of the edge heat-kernel/Casimir law itself.
That remains the compact-paper task.

\subsection{Direct feed: \texttt{papers.reality.a.01}}

The consensus paper needs a regulated patch net with explicit local state spaces, shared
observables, and repair maps. This note now supplies:
\begin{enumerate}[leftmargin=2em]
\item a concrete patch cover on a finite screen;
\item explicit overlap observable families and mismatch syndromes;
\item a finite local repair interface $\mathcal R_{PQ}$ and asynchronous update loop.
\end{enumerate}
What remains open are the theorem-level claims about termination, local confluence, and gauge
covariance of those repairs.

\subsection{Direct feeds: \texttt{papers.observers.b.03} and \texttt{papers.observers.e.18}}

The observer-side tasks need a concrete record layer and explicit checkpoint-like data. This note
now supplies:
\begin{enumerate}[leftmargin=2em]
\item a record subalgebra candidate built from robust pointer observables;
\item an operational observer criterion
$O_{\mathrm{ref}}=(P_O,\mathcal A(P_O),\rho_O,R_O,\mathcal U_O)$;
\item explicit shared record summaries and update interfaces that later continuation or
measurement theorems could quantify.
\end{enumerate}
What remains open are the theorem-level Born-rule package, identity criteria, restoration maps,
and accumulated-error bounds.

\subsection{Likely later feeds}

This note is likely to feed later compact quantitative tasks, later repair-law and confluence
tasks, and later quantum-record / continuum-lift tasks, but only indirectly. It gives those
tasks a concrete regulated surface on which quantities can be defined and measured. It does not
claim to close their later derivation packages in phase 1.

So the correct role of this note inside the completion tree is narrow but important. It makes the
fixed-cutoff build surface explicit without replacing the later derivation or theorem packages.

\section{Conclusion}

The main missing bridge for a constructive OPH microphysics program is no longer hard to say.
One must go from finite local screen registers to gauge-invariant patch algebras, from patch
algebras to overlap observables and edge sectors, from those shared observables to stable
records, from records to observer patterns, and from observer patterns to gauge-covariant local
synchronization rules. This paper now spells that bridge out at the level of a finite reference
architecture: explicit shared observables on overlaps, an explicit mismatch syndrome and repair
interface, an octahedral $\mathbb Z_2$ pilot, an $S_3$ nonabelian upgrade path, and a validation
package centered on gauge violation, overlap mismatch, record lifetime, and schedule divergence.

If the reference architecture fails, that is useful: it tells us that this class of finite
screen models does not realize the needed OPH interfaces. If it succeeds, that still does not
identify the final UV completion or prove the continuum branch. What it does provide is the
phase-1 architecture needed before stronger claims can be asked responsibly. The correct reading
rule is therefore simple: this note is a finite-cutoff build surface for later OPH completion
tasks, not the claim that those tasks are already done.

\end{document}
